% MẪU 1 — Chuyên đề cơ bản: chương gồm nhiều bài.
% Mỗi bài: A. Lý thuyết, B. Ví dụ (tự luận), C. Bài tập trắc nghiệm chia theo dạng;
% trong mỗi dạng lần lượt: trắc nghiệm, đúng sai, trả lời ngắn.
% Đổi lớp: chỉ sửa dòng \usepackage{toan10} thành toan11, toan12...
\documentclass[12pt]{article}
\usepackage{toan10}
\begin{document}
\loai{1}\chuong{1}

% ===================== BÀI 1 =====================
\tieude{\tentailieu}{Bài 1. Mệnh đề}

\phan{A. Lý thuyết}
\begin{kienthuc}
\textbf{1. Mệnh đề.} Mệnh đề là một câu khẳng định đúng hoặc sai. Một mệnh đề không thể vừa đúng vừa sai.

\textbf{2. Mệnh đề phủ định.} Phủ định của mệnh đề $P$ kí hiệu là $\overline{P}$.
\end{kienthuc}

\phan{B. Ví dụ}
\begin{vd}
Trong các câu sau, câu nào là mệnh đề? Xét tính đúng sai của mệnh đề đó.
\begin{cacy}(2)
\item $2+3=5$.
\item Hôm nay trời đẹp quá!
\end{cacy}
\giai{a) Là mệnh đề đúng. b) Là câu cảm thán, không phải mệnh đề.}
\end{vd}

\phan{C. Bài tập trắc nghiệm}
\dang{Nhận biết mệnh đề, xét tính đúng sai}
\muctn
\begin{tn}[NB]
Câu nào sau đây là mệnh đề?
\pa{Mấy giờ rồi?}{\dung{$3>1$.}}{Bạn học lớp mấy?}{Đẹp quá!}
\loigiai{$3>1$ là câu khẳng định đúng.}
\end{tn}
\begin{tn}[TH]
Mệnh đề nào sau đây sai?
\paa{$\sqrt{4}=2$.}{$2$ là số nguyên tố.}{\dung{$\pi$ là số hữu tỉ.}}{$5>-1$.}
\end{tn}

\mucds
\begin{ds}[TH]
Cho mệnh đề $P$: ``$x^2\ge 0$ với mọi số thực $x$''.
\ydung{$P$ là mệnh đề đúng.}
\ysai{Mệnh đề phủ định của $P$ là ``$x^2<0$ với mọi số thực $x$''.}
\ydung{Mệnh đề phủ định của $P$ là ``Tồn tại số thực $x$ để $x^2<0$''.}
\ysai{Phủ định của $P$ là mệnh đề đúng.}
\loigiai{Phủ định của ``với mọi'' là ``tồn tại''.}
\end{ds}

\muctln
\begin{tln}[VD]
Có bao nhiêu số nguyên dương $n<10$ để mệnh đề ``$n$ chia hết cho 3'' đúng?
\traloi{3}
\loigiai{Đó là $3$, $6$, $9$.}
\end{tln}

\dang{Mệnh đề kéo theo, mệnh đề đảo}
\muctn
\begin{tn}
Mệnh đề đảo của ``Nếu $a=b$ thì $a^2=b^2$'' là
\paaa{Nếu $a\ne b$ thì $a^2\ne b^2$.}{\dung{Nếu $a^2=b^2$ thì $a=b$.}}{Nếu $a^2\ne b^2$ thì $a\ne b$.}{$a=b$ khi và chỉ khi $a^2=b^2$.}
\end{tn}

\bangdapan   % bảng đáp án của bài (chỉ bản GV)

% ===================== BÀI 2 =====================
\newpage\baimoi
\tieude{}{Bài 2. Tập hợp và các phép toán trên tập hợp}
\phan{A. Lý thuyết}
\begin{kienthuc}...\end{kienthuc}
\phan{B. Ví dụ}
\phan{C. Bài tập trắc nghiệm}
\dang{...}
\end{document}
